\chapter{Literature Survey and Research Gap}
Earlier work has shown that quantization and the choice of serving engine can each change inference energy, and that energy-aware routing between models can be effective. What it leaves open are three narrower questions: how quantization implementations compare when paired directly, whether energy measurement stays valid under host interference, and how prompt-level precision routing performs when evaluated in joules.

\section{Review method}
Sources were located through arXiv, ACL Anthology, USENIX, Crossref and targeted web searches, in November 2026. Search terms combined LLM inference energy, quantization, FP8, AWQ, bitsandbytes, vLLM, precision routing and energy-aware routing. Papers were retained if they measured or modelled inference energy, implemented low-precision execution, or routed queries using cost or quality signals. The search was targeted, not a systematic review, so the absence of a work from this chapter does not establish that none exists.

\section{Evolution of the problem}
The first papers in this area established the cost of training and argued for efficiency as an evaluation criterion~\cite{b1,b2,b3,b4}. Because they relied largely on estimation and concern training, they say little about the per-token cost of serving. Samsi et al. shifted attention to inference, benchmarking LLaMA on V100 and A100 GPUs at multi-GPU scale~\cite{b5}.

Later studies broke that cost down. Vellaisamy et al. model a request as a fixed cost plus a marginal per-token cost, and show on H100 and H200 that token energy can fall while request energy rises~\cite{b8}. Wilkins et al. fit workload-based energy models for offline scheduling~\cite{b9}. At cluster level, Stojkovic et al. and DynamoLLM adjust settings such as instance count, parallelism and GPU frequency, with DynamoLLM reporting 53\% energy savings at service level~\cite{b10,b11}, and POLCA addresses power oversubscription~\cite{b12}. In these systems the model's numeric precision is fixed or peripheral, which leaves it as a separate question.

\section{Energy of quantized inference}
At first sight the published findings disagree, but differences in baseline and workload account for much of the disagreement.

Poddar et al. report that bitsandbytes 8-bit and 4-bit raise energy to almost 2 times at equal batch size, because computation still runs in 16-bit, and that savings arise from the larger batches quantization permits~\cite{b6}. Their software tools, CodeCarbon and Carbontracker, sampled every 10 s in the main runs, which suits long workloads better than individual generations.

Basta evaluates ten weight-only strategies in Transformers on four GPUs, measuring GPU energy with Zeus, and reports that bitsandbytes int8 and EETQ int8 save up to 4 times against FP32 on short generation~\cite{b47}. This does not directly contradict Poddar et al., because Basta's baseline is FP32 and not a 16-bit format, the cited case is prefill-heavy (1,000 input tokens and 1 output token), and dequantization overhead is explained as amortized during prefill. To isolate quantization, Basta deliberately avoids vLLM and TensorRT-LLM, which leaves engine-level kernels untested, and lists the separation of memory-access from compute energy as future work.

TokenPowerBench reports that FP8 cuts energy per token by about 30\% on Llama 3 405B, and that vLLM and TensorRT-LLM reduce energy per token by 25-40\% against Transformers~\cite{b7}. It treats quantization and engine as separate sensitivity studies at a scale far above a single-GPU deployment, and reports neither repeat-run variance nor sensor validation.

Taken together, these studies suggest that the answer depends on baseline, workload phase and implementation. None of the works reviewed, however, provides a paired comparison of bitsandbytes and native vLLM tiers against float16 of the same model on one GPU model in a decode-heavy workload.

\section{Validity of energy measurement}
The instruments have limits of their own. Yang et al. tested NVIDIA's on-board sensor on over 70 GPUs against an external meter and found that only 25\% of runtime is sampled for power on A100 and H100~\cite{b17}. Zeus, LLMCarbon, Carbontracker and Henderson et al. made measurement and reporting practical~\cite{b13,b14,b15,b16}, but being easy to measure says nothing about whether a given run was undisturbed.

Isolation practices vary. Basta verifies an otherwise idle GPU with nvidia-smi~\cite{b47}, while Vellaisamy et al., working on a shared research cluster, use warm-up runs and accept replay windows only within 5\% variation, without describing co-tenant detection~\cite{b8}. Both approaches focus on the GPU, and neither addresses a case in which the GPU is unshared yet the workload is slowed from the host side. The systems literature has long warned that innocuous setup differences can bias conclusions~\cite{b18} and has set rules for reporting performance~\cite{b19}, but these warnings have not been applied, in the works reviewed, to joules per token on shared nodes.

\section{Dynamic precision}
LLM.int8() and NF4 made low-bit weights accessible~\cite{b20,b21}, and GPTQ, AWQ and SmoothQuant improved post-training accuracy~\cite{b22,b23,b24}, but each of these methods selects one static precision.

Recent work varies precision at run time. Any-Precision LLM overlays several bit-widths in one memory footprint, but reports latency, throughput and memory, and does not specify a selection policy~\cite{b25}. DP-LLM assigns precision per layer, on the observation that layer sensitivity changes across decoding steps~\cite{b50}. FlexQuant switches precision during token generation using entropy and KL divergence and reports a 1.3 times speedup~\cite{b51}, and MoBiQuant, which uses a token-aware router, reports throughput gains of up to 1.34 times~\cite{b52}.

These methods need custom kernels or engines and judge success by speed, not joules. Their signals are also internal to the model and act at layer or token level. The premise they share, that sensitivity varies within a generation, offers a possible explanation for why a single decision made from the prompt text might be too coarse, and this thesis tests that prompt-level decision directly.

\section{Query routing}
A recent survey classifies routers by when they decide, what information they use and how they compute the decision, and lists text length and syntactic complexity among heuristic difficulty signals~\cite{b53}. It presents these as lightweight options and provides no systematic empirical comparison with learned routers.

FrugalGPT cascades across paid APIs~\cite{b28}. RouteLLM learns from preference data, and its results caution against assuming a router works: Arena-only routers perform at or near random on MMLU and below random on GSM8K, and they beat random only with augmented data~\cite{b29}. Hybrid LLM trains a BERT-style encoder on the quality gap between a small and a large model, and measures cost as the share of queries sent to the small model~\cite{b30}. AutoMix verifies a small model's output and routes with a POMDP~\cite{b31}, and RouterBench provides an evaluation benchmark~\cite{b32}. These routers target price or call share, and each estimates quality in some form, not surface complexity alone.

Energy-aware routing has also been studied. GreenServ routes among 16 models using a contextual bandit, measuring energy with Zeus and reporting 31\% less energy than a random baseline with higher accuracy~\cite{b48}. In its ablation, task type gave the largest reduction in regret, while the text complexity feature slightly increased regret. Siddiqui et al. train a language-model router on queries profiled offline, polling NVML at 200 ms on A100 GPUs, and report about 23\% less energy than UniRoute at nearly the same accuracy~\cite{b49}. Both route among different models, and neither lists quantized variants in its pool. GreenServ's Random baseline is not described as matched to the router's mix, so routing intelligence cannot be separated from a shift in model mix.

Fuzzy logic~\cite{b33,b34} has an established rationale for interpretable control. Readability inputs such as the Flesch-Kincaid grade~\cite{b35}, however, were designed for human readers and not for predicting quantization sensitivity, so that link requires empirical testing.

\section{Comparison of existing approaches}

\begingroup\footnotesize\setlength{\tabcolsep}{4pt}
\begin{longtable}{>{\raggedright\arraybackslash}p{30.1mm}>{\raggedright\arraybackslash}p{24.8mm}>{\raggedright\arraybackslash}p{26.5mm}>{\raggedright\arraybackslash}p{23.6mm}>{\raggedright\arraybackslash}p{31.0mm}}
\toprule
\textbf{Work} & \textbf{Variable studied} & \textbf{Energy evidence} & \textbf{Limitation for this thesis's question} & \textbf{Relationship to this thesis} \\
\midrule
\endhead
Poddar et al.~\cite{b6} & bitsandbytes, batch size, tasks & CodeCarbon, Carbontracker; about 10 s sampling & Single quantization path; software tools & Stack A tests a similar path at batch 1 with a counter \\*
Basta~\cite{b47} & Ten libraries, four GPUs & Zeus; GPU only & FP32 baseline; no serving engine & Different baseline and workload; complementary \\
TokenPowerBench~\cite{b7} & Models, engines, FP8 & NVML/DCGM, RAPL, IPMI & FP8 at 405B only; validity checks not stated & Stack B finds comparable direction at 1B and 8B \\
Vellaisamy et al.~\cite{b8} & Request vs token energy & NVML cumulative counter & BFloat16 only; co-tenant control not described & Same instrument; supplies the per-request caution \\
Yang et al.~\cite{b17} & Sensor accuracy & External meter comparison & Sensor, not host interference & Basis for a stated threat to validity \\
Any-Precision LLM~\cite{b25} & Overlaid bit-widths & Not measured & No selection policy & Same idea of tiers; this thesis adds joules \\
DP-LLM, FlexQuant, MoBiQuant~\cite{b50,b51,b52} & Layer or token precision & Not the reported metric & Custom execution; speed-based evaluation & Finer granularity; motivates the coarseness hypothesis \\*
RouteLLM, Hybrid LLM~\cite{b29,b30} & Router between two models & Not measured & Cost as calls or price & Provides random baseline and quality-prediction ideas \\*
GreenServ, Siddiqui et al.~\cite{b48,b49} & Energy-aware routing among models & Zeus; NVML polled at 200 ms & Different models, not precision tiers & Closest routing work; agrees that text complexity is weak \\*
\bottomrule\noalign{\nobreak}
\end{longtable}
\addtocounter{table}{-1}
\endgroup

\section{Research gap}
The following gaps are claimed for the works reviewed.

\begin{itemize}
\item \textbf{G1, paired implementation comparison.} Quantization energy has been studied within Transformers~\cite{b6,b47} and, separately, at very large scale across engines~\cite{b7}. No reviewed work compares bitsandbytes and native vLLM tiers against float16 of the same model, on one GPU model, in a decode-heavy workload at batch sizes 1 and 8, while also reporting the accuracy trade-off.
\item \textbf{G2, validity under host interference.} Existing checks focus on the sensor~\cite{b17} or on GPU occupancy~\cite{b47}. No reviewed work shows that repeat-run agreement, a co-tenancy check and an automated validator can together accept an energy reading inflated several times by a slowdown not visible at those checks.
\item \textbf{G3, prompt-level precision routing in joules.} Dynamic-precision work reports speed at layer or token level~\cite{b25,b50,b51,b52}, and energy-aware routing selects among different models~\cite{b48,b49}. No reviewed work evaluates routing between precision tiers using measured energy and a random baseline matched to the router's mix, or tests surface features as direct predictors of sensitivity to quantization.
\end{itemize}
